We compare exploration schemes for tabular Q-learning on chain and multi-room gridworlds with one sparse goal reward: $\epsilon$-greedy, count-based bonuses and an RND-style bonus (a small numpy predictor distilling a fixed random target), all reimplemented, with controls for the constant per-step penalty (an optimistic offset) that the bonus systems carry, and a noisy-TV cell that emits a random token. This is a small CPU study: five seeds per cell, one fixed hyperparameter setting. The count bonus found the reward in every run on all eight tasks, while $\epsilon$-greedy found it in one of five seeds on chain\_20 and in none on chain\_40, room\_4 and room\_6. This gap is not evidence for the count signal: a control with the penalty alone and no bonus is within one standard deviation of the count bonus on every noise-free task (uncorrected Welch $p\geq0.288$), and the count bonus without the penalty fails on the larger tasks. The RND-style bonus found the reward in 39 of 40 runs but was slower than both the count bonus and the penalty-only control on every task (18435 steps against 11476 and 11921 on room\_4), so here its normalised novelty signal slowed exploration relative to no bonus. We document a failure case: the first version of our normaliser divided by a running standard deviation that is zero when the first errors coincide, giving bonuses up to $1.44\times10^{7}$; 19 of the 20 runs with such a spike never found the reward, all 20 runs without one did, and either a warm-up or a clip removes the failure. RND is designed to be robust to noisy-TV stochasticity, and consistent with that the TV delayed neither the RND nor the observation-count bonus; their share of steps on the TV rose with the number of tokens (0.016 to 0.029 for RND on room\_4\_tv). These are single-implementation, small-scale results, not a verdict on RND at scale.
